\subsection{Conversion can change another expense and depend on another instrument}
\label{dependent-eps-consequences}

Reallocating a fixed amount of earnings is only one way that assumed conversion
can change the numerator. Removing interest expense may increase the earnings
on which employee remuneration depends. IAS~33.35 expressly requires the
resulting change in nondiscretionary profit-sharing expense to enter the EPS
calculation, together with the related tax effects under paragraphs 33--34
\citep[IAS~33.33--35]{eu20231803}. Calculating each conversion once in isolation
can therefore miss a consequence of converting two instruments together.

Suppose profit before a contractual employee bonus is 90. The bonus is half of
profit above 100, and complete tax expense is one fifth of profit after the
bonus. These are supplied terms for the example. There are 100 ordinary shares,
with no participation rights or other attribution; continuing and total
earnings coincide. The reported bonus is zero, tax is 18 and earnings are 72.
Each of two full-year convertible debt series would remove interest expense of
20 and add 20 ordinary shares. Assume all other conversion consequences have
been accounted for.

Converting either series alone raises profit before the bonus to 110. The bonus
becomes 5, tax becomes 21 and earnings become 84. The complete after-tax increase
is 12, and the hypothetical EPS is \(84/120=7/10\). Converting both raises
profit before the bonus to 130. Now the bonus is 15, tax is 23 and earnings
are 92, giving hypothetical EPS \(92/140=23/35\). Adding the two isolated
increases would instead give \(72+12+12=96\). That four-unit overstatement
comes from failing to recalculate the common bonus and its tax consequence.

For a set \(S\) of assumed conversions, let \(I_i\) denote interest saved by
series \(i\), and \(b_i\) its incremental ordinary shares. Write \(P\) for
profit before the bonus, \(H\) for its threshold, \(q\) for the bonus rate
and \(t\) for the complete proportional tax rate. The supplied terms imply
\[
 \begin{aligned}
 P(S)&=P+\sum_{i\in S}I_i,&
 B(S)&=q\max(P(S)-H,0),\\
 T(S)&=t\{P(S)-B(S)\},&
 N(S)&=P(S)-B(S)-T(S),\\
 D(S)&=n+\sum_{i\in S}b_i,&
 r(S)&=N(S)/D(S).
 \end{aligned}
\]
Thus the first conversion adds 12 to earnings, while the second adds only
\(92-84=8\). With \(q=0\), the bonus disappears and the two after-tax
interest effects become additive. The formula describes this remuneration and
tax model; it does not decide whether a real obligation meets recognition
requirements or establish its tax treatment. The loss illustrations use zero
tax rather than assuming that every loss creates a recoverable tax asset.

A complete hypothetical calculation still needs the applicable dilution rule.
For example, change only initial profit before the bonus to 80. Basic earnings
are 64 and EPS is \(16/25\). Either conversion alone gives earnings 80 and
EPS \(2/3\), an increase. Both together give earnings 88 and EPS \(22/35\),
below basic EPS. Merely taking the smallest of these values would bypass the
individual-series requirements of IAS~33.41--44
\citep[IAS~33.41--44]{eu20231803}. The
\texttt{ConsequentialEPSWorlds/1} method therefore reports the complete four
worlds and leaves final dilution selection unevaluated. The earlier additive
ordering proof does not apply to these joint earnings effects.

There is a separate, explicit dependency for contingently issuable potential
shares. IAS~33.57 first tests whether the instruments may be assumed issuable
under paragraphs 52--56, then applies their instrument-specific dilution rules.
Its final sentence also requires assumed exercise or conversion of similar
outstanding potential shares that are not contingently issuable
\citep[IAS~33.52--57]{eu20231803}. A satisfied contingency alone is insufficient.

Consider profit 100 and 100 ordinary shares. A similar outstanding option
series contains 20 options at strike 6; a contingent option series contains
20 at strike 5. Average market price is 10. For this conditional example,
similarity has been separately admitted; different exercise prices neither
establish nor automatically disprove that finding. Both agreements cover the
full reporting period, with no future-service proceeds, earnings effects,
participation or additional dependencies. The contingent option is excluded
from basic shares. Its nonfinancial condition requires five stores, assessed
at the reporting-period end while the contingency period remains open.

Under paragraphs 52 and 56, use the period-end condition status. Five observed
stores qualify; a forecast of a fifth opening next year does not. The
outstanding option supplies \(20-20(6)/10=8\) incremental ordinary shares.
Its exercise is dilutive, reducing EPS to \(100/108=25/27\). With the
condition satisfied, the contingent option supplies another
\(20-20(5)/10=10\) shares, giving diluted EPS \(100/118=50/59\).
Both have zero earnings increments, so testing the prerequisite first respects
their tied incremental-earnings ratios while enforcing paragraph 57.

Now give the similar outstanding option strike 12, retaining the example's
explicit similarity finding. It is out of the money at average price 10 and
its exercise is not assumed. The contingent option would reduce EPS to
\(100/110=10/11\) in isolation, but paragraph 57's prerequisite is unsatisfied.
Both are excluded and EPS remains 1. If instead only four stores are observed,
the original outstanding option remains dilutive but the contingent option is
excluded, leaving EPS \(25/27\). At zero or negative earnings neither
zero-earnings-effect option strictly reduces EPS or increases loss per share.

The \texttt{ContingentOptionEPS/1} method records both tests and each exclusion
reason. Missing similarity or complete terms leaves the answer unresolved;
future observations, incompatible periods and unsupported conditions cannot
produce a partial result. Seven consequence cases and eight option cases check
these distinct calculations. Arbitrary linked conversions and a general
selection method for nonadditive earnings remain separate source and derivation
requirements, with the complete conditional worlds preserved for that work.
